\section{Conclusion}
\label{sec:conclusion}

This paper presents \softrobotagent{}, a physics-grounded research framework for autonomous soft-robot design and control. Its central idea is to connect agent-directed physical hypotheses with task-conditioned mathematical problems, scoped sensitivity analysis, and independent physical evaluation. The M0--M3 semantics distinguish mechanism, fixed-input response, re-optimized task capability, and closed-loop performance, while the parameterized research interface allows active and passive robots to use different inner physical solvers without changing the overall research organization.

The completed task portfolio and the nine-segment PHS-2D case provide evidence of multi-task execution, adaptive research decisions, passive structural design, and local sensitivity consistency. The ongoing matched A/B/C comparisons will quantify the additional value of direct mathematical tools and the complete research architecture. Physical-robot demonstrations will provide a separate validation layer using the project's existing hardware records.

The broader objective is to make mathematical and physical evidence actionable within an autonomous research process: not only to generate a candidate, but to determine what should be investigated, why a design should change, and how the resulting change should be verified.
